% Standalone appendix document - compiles on its own:
%   latexmk -pdf -outdir=build 05-agenda-2026-10-07.tex
\documentclass[a4paper,11pt]{article}
\usepackage[danish]{babel}
\usepackage{../../appendix-style}
\usepackage{pgfplotstable}
\pgfplotsset{compat=1.18}

\title{12.5 Dagsorden -- 07-10}
\author{SW4PRJ4 -- Group 3}
\date{07-10-2026}

\begin{document}
\maketitle

\section*{Mødeinformation}

\textbf{Dato:} 07-10-2026 (onsdag)

\textbf{Tid:} 10:00

\textbf{Sted:} Incuba

\section*{Dagsordenspunkter}

\begin{enumerate}
    \item \textbf{Opsamling fra vejledermødet 6/10} (ca. 15 min)
    \begin{itemize}
        \item Referatet er bilag 13.4. Til stede var Alexander H, Alexander D, Frederik, Ida D og Ida F.
        \item \textbf{Kilder:} Vi må kun henvise til originale kilder (dokumentation og litteratur), ikke til slides på Brightspace. Det rammer 29 henvisninger i den tekniske analyse og i bilag 10.3 (PMU-81).
        \item \textbf{Diagrammer:} Vi følger Projekt 4-prototypen og laver diagrammerne løbende i nogenlunde UML. Michel lægger mere vægt på, at det virker, og at vi forstår det, end på det formelle.
        \item \textbf{Auth-flow:} Hele flowet fra klient til Keycloak og videre til backend skal beskrives med et diagram, hvor der står teknisk tekst på pilene. Kunde-appen og admin-panelet kan have forskellige flows (PMU-77).
        \item Credentials går kun mellem klient og Keycloak. Backend ser kun tokens.
        \item \textbf{Claims og roller:} Vi skal beslutte, hvilke claims de forskellige brugere får, og hvordan man bliver admin (PMU-75).
        \item \textbf{Flows mellem noderne:} Vi skal løbende dokumentere, hvordan data flyder mellem kunde-app, admin-panel, backend og database, så rapporten viser, at vi har forstået applikationen (PMU-78).
        \item \textbf{Keycloak:} Rapporten skal sige, hvilke features vi bruger, og hvilke vi har fravalgt (PMU-79). Vi skal forhindre zombie-brugere, nu hvor Keycloak ligger offentligt (PMU-80).
        \item \textbf{Expo:} Rapporten skal beskrive, at kunde-appen udgives gennem Expo og ikke direkte til Google Play (PMU-82).
        \item Michels bud på næste skridt: onboarding af en bruger og helt styr på Keycloak-flowet, også admin-login i admin-panelet.
    \end{itemize}
    \item \textbf{Sprint review -- sprint 2} (ca. 45 min)
    \begin{itemize}
        \item Sprint goal: gøre sprint 1 færdig og levere et E/R-diagram og et low fidelity UI-mockup. Blev det nået?
        \item Demo: hver af os gennemgår det, vi selv har lavet i sprinten.
        \item I \emph{Done}: 9 kort (25 point). PMU-2, PMU-3, PMU-4, PMU-12, PMU-13, PMU-18, PMU-19, PMU-71 og PMU-72.
        \item Ikke færdigt: PMU-70 (login i kunde-appen med Keycloak, 3 point) er i \emph{In Progress}.
        \item Velocity for sprint 2: committed 28 point, done 25 point, spillover 3 point, scope 62 point. PMU-71 (2 point) og PMU-72 (3 point) fik først point-label ved sprintens afslutning og er talt med i committed. Tallene står i \texttt{docs/assets/data/sprint-velocity.csv}.
        \item Vi blev færdige med 60 \% af det, vi tog ind i sprint 1, og 89 \% i sprint 2. Figur~\ref{fig:velocity} viser de to sprints ved siden af hinanden.
    \end{itemize}

    \begin{figure}[H]
        \centering
        \pgfplotstableread[col sep=comma]{../../../assets/data/sprint-velocity.csv}\velocitytable
        \begin{tikzpicture}
            \begin{axis}[
                ybar, bar width=14pt,
                width=0.8\textwidth, height=6cm,
                xlabel={Sprint}, ylabel={Point},
                xmin=0.5, xmax=6.5, xtick={1,...,6}, ymin=0,
                legend pos=north east,
                nodes near coords, nodes near coords style={font=\footnotesize},
            ]
                \addplot[fill=blue!35, draw=blue] table[x=sprint, y=committed]{\velocitytable};
                \addlegendentry{Committed}
                \addplot[fill=red!45, draw=red] table[x=sprint, y=done]{\velocitytable};
                \addlegendentry{Done}
            \end{axis}
        \end{tikzpicture}
        \caption{Velocity pr. sprint. Afstanden mellem committed og done er faldet fra 23 point i sprint 1 til 3 point i sprint 2: vi tog mindre ind og blev færdige med næsten det hele.}
        \label{fig:velocity}
    \end{figure}
    \item \textbf{Retrospective} (ca. 30 min)
    \begin{itemize}
        \item Opfølgning på sidste retrospective: har demoen givet overblik over den anden undergruppes arbejde? Bliver E/R-diagrammet brugt som overbliksdiagram?
        \item Hvad gik godt?
        \item Hvad gik skidt?
        \item Hvad ændrer vi i sprint 3? Hver handling bliver et kort i Kaneo.
    \end{itemize}
    \item \textbf{Sprint planning -- sprint 3 (7/10 -- 27/10)} (ca. 60 min)
    \begin{itemize}
        \item Forslag: sprint 3 forlænges en uge til 27/10 på grund af efterårsferien i uge 42. Sprint 4 -- 6 rykkes en uge, så sprint 6 slutter 8/12, tre dage før aflevering.
        \item Kapacitet: hvor mange timer har hver af os i perioden?
        \item Spillover: kort, der ikke er færdige, beholder label \emph{Sprint 2} og får \emph{Sprint 3} oveni.
        \item Efter vejledermødet er der oprettet ni kort, PMU-74 -- PMU-82. De ligger allerede i \emph{To Do} med label \emph{Sprint 3}, start 7/10 og frist 27/10, men er ikke ready: de mangler point-label og ejer. Gruppen skal godkende, at de er med i sprinten.
        \begin{itemize}
            \item Kode: PMU-74 onboarding af ny bruger, PMU-75 claims og roller i Keycloak, PMU-76 admin-login i admin-panelet, PMU-80 zombie-brugere.
            \item Rapport: PMU-77 auth-flow, PMU-78 flowdiagram for gem-flow, PMU-79 Keycloak-features, PMU-81 originale kilder, PMU-82 udgivelse gennem Expo.
            \item Forslag til estimater står i beskrivelsen på hvert kort og giver 26 point i alt uden PMU-80, som først kan estimeres, når den er afklaret. Med PMU-70 er det 29 point, før de gamle kandidater regnes med.
        \end{itemize}
        \item Fra retrospectiven: tre nye kort ligger i backloggen uden point-label og ejer. PMU-83 test først, PMU-84 demo-video på Drive og PMU-85 gruppesang. Skal de med i sprint 3, og skal de have point?
        \item Test først (PMU-83): for hvert kodekort, der trækkes ind, skriver vi målet og testtilfældene på kortet, før implementeringen starter. Skal det være en del af ready?
        \item PMU-74 har krav FR-01 ligesom PMU-70. Hvor går grænsen mellem login og onboarding, og hvad består onboarding af efter kravspecifikation 1.4?
        \item PMU-27 er slået sammen med PMU-70, som har overtaget krav FR-01 og definition of done. PMU-27 skal slettes, og kravlinjen skal ind i beskrivelsen på PMU-70.
        \item PMU-30 (GUI-skitser) overlapper med mockuppet i PMU-72. Er PMU-30 stadig nødvendig?
        \item PMU-73 (kravspec 1.4) ligger i backloggen, men har label \emph{Sprint 2}. Skal den med i sprint 3, eller skal labelen fjernes?
        \item Kandidater til sprint 3: husstandsslicen PMU-24 -- PMU-26 og PMU-29, konto PMU-28 og rapportkortene PMU-31 -- PMU-33.
        \item Planning poker på de kort, der trækkes ind. Ingen af dem har point-label endnu.
        \item Ready-kort fra backlog til \emph{To Do}: \texttt{Krav: FR-xx} som første linje, definition of done, point-label, epic-label og ejer.
        \item Sprint goal (én sætning, der nævner de krav sprinten leverer).
        \item Ejer og undergruppe på hvert kort.
        \item Point i \emph{To Do} (committed).
    \end{itemize}
    \item \textbf{Eventuelt}
\end{enumerate}

\end{document}
